\uuid{YcAu}
\titre{Implication/équivalence}
\auteur{bodin ; hanani ; mzari}
\organisation{exo7}
\licence{CC BY-NC-SA 4.0}
\datecreate{2020-01-20}
\niveau{L1}
\module{Algèbre}
\chapitre{Logique, ensemble, raisonnement}
\sousChapitre{Absurde et contraposée}
\difficulte{2}
\theme{Raisonnement, implication, équivalence}
\exo7qcmid{lille-1.logique.023}

\contenu{
\begin{question}
\qclassification{100.03,100.04}
\qtags{motcle=implication/équivalence}
J'effectue le raisonnement suivant avec deux fonctions $f,g : \Rr \to \Rr$.
$$\forall x \in \Rr \quad f(x)\times g(x) = 0$$
$$\implies \forall x \in \Rr \quad \big(f(x) = 0 \text{ ou } g(x) = 0\big)$$
$$\implies \big(\forall x \in \Rr \quad f(x) = 0\big) \ \text{ ou } \ \big(\forall x \in \Rr \quad g(x) = 0\big)$$
\begin{answers}
    \bad{Ce raisonnement est valide.}

    \bad{Ce raisonnement est faux car la première implication est fausse.}

    \good{Ce raisonnement est faux car la seconde implication est fausse.}

    \bad{Ce raisonnement est faux car la première et la seconde implication sont fausses.}
\end{answers}
\begin{explanations}
On ne peut pas distribuer un "pour tout" avec un "ou".
\end{explanations}
\end{question}
}
